\section{Attribute Synthesis: Algorithm and Guarantees}
\label{sec:mathematical-synthesis}
We now construct RHS attributes from the LHS axis representation. Logical
axes and the target output order $F$ always come from the LHS. The algorithm
has three passes: find the axes that survive at each node, check whether
the required orders are possible, and construct attributes. It does not
enumerate possible attribute assignments.

\subsection{Canonical Intermediate Axes}
\label{sec:canonical-axes}
The question is how to decide which axes each intermediate RHS result must
keep. An axis must remain if it is needed in the final output or must connect
to another input later. Otherwise, we remove it as soon as the operation
allows.

We first identify the axes that must be kept, called the \emph{mandatory
axes}. For a node $v$, let $U_v$ be the inputs in its subtree, and let
$\sigma(a)$ be the inputs that contain logical axis $a$. The mandatory
set $M_v$ contains axes used in this subtree that are also needed in the
final output $F$ or in an input outside the subtree:
\begin{equation}
 M_v=\left\{a\in\mathcal A\;\middle|\;
 \sigma(a)\cap U_v\ne\varnothing
 \land (a\in\operatorname{set}(F)\lor\sigma(a)\nsubseteq U_v)
 \right\}.
 \label{eq:mandatory-axes}
\end{equation}
An operation may also have to keep an axis that is no longer needed.
For example, a contraction cannot sum out an axis present in only one
operand; it must keep that axis as a free axis.

Let $P_v$ be the axes that our algorithm keeps in the output of node
$v$.\footnote{$P_v$ is a set of logical axes; $P$ in
Section~\ref{sec:axis-model} is the set of input ports.}
This set includes every mandatory axis, so $M_v\subseteq P_v$.
We call it the \emph{canonical live set} because the algorithm follows one
fixed rule: remove each unneeded axis at the first operation that can
legally remove it.

We compute these sets from the RHS inputs toward the root, using
Table~\ref{tab:canonical-live}. Here $I_X$ lists the axes of input $X$ in
their original order, using the logical axes identified on the LHS.
\begin{table}[!htbp]
\caption{Axes kept at each RHS node. A unary node has child $u$.}
\label{tab:canonical-live}
\centering
\begin{tabular}{@{}lll@{}}\toprule
Node & Axes kept ($P_v$) & Reason\\\midrule
Input $X$ & $\operatorname{set}(I_X)$ & Start with all input axes.\\
Permutation & $P_u$ & Keep all axes and change their order.\\
Unary summation & $P_u\cap M_v$ & Keep only mandatory axes.\\
Contraction & $B_v\cup LF_v\cup RF_v$ & Keep batch and free axes.\\\bottomrule
\end{tabular}
\end{table}
For a contraction, a shared axis is kept as a batch axis if it is mandatory;
otherwise, it is summed out. Axes present in only one operand remain free.
For children $\ell,r$, the batch, contraction, left-free, and right-free
axis sets are therefore
\begin{equation}
B_v=P_\ell\cap P_r\cap M_v,\quad
K_v=(P_\ell\cap P_r)\setminus M_v,\quad
LF_v=P_\ell\setminus P_r,\quad
RF_v=P_r\setminus P_\ell.
\label{eq:canonical-roles}
\end{equation}
An axis is \emph{complete} in a subtree when all inputs containing it are
inside that subtree. Only complete axes absent from the final output can
be removed. Contraction removes such axes when they are shared; unary
summation can remove any such axis in its input.

Applying these rules from the inputs upward keeps every mandatory axis,
and every retained axis comes from an input in the subtree. At the root,
the remaining axes must be exactly those required by the LHS output:
\begin{equation}
 P_{\mathrm{root}}=\operatorname{set}(F).
 \label{eq:root-live}
\end{equation}
Extra axes may survive intermediate nodes, but none may remain in the final
output. Appendix~\ref{app:canonical-axes-example} provides a concrete example
of how these axis sets are computed for batched matrix multiplication
followed by a row sum, including a case where $M_v\subsetneq P_v$.

\subsection{Why the Canonical Choice Does Not Lose Solutions}
\label{sec:canonicalization}
Let $\operatorname{Ord}_v(\beta)$ be the ordered list of logical axes
in node $v$'s output when the RHS attributes are assigned by $\beta$.
The notation $S|_Q$ denotes the subsequence obtained by removing from $S$
every element not in $Q$, while preserving the order of the remaining elements. By
Lemma~\ref{lem:axis-exactness}, any valid RHS has the LHS partition, so its
intermediate axes can be named using LHS classes.

\begin{lemma}[Canonicalization preserves feasibility]
\label{lem:canonicalization}
If a valid assignment $\beta$ exists, a valid assignment $\beta'$ exists with
\begin{equation}
 \forall v,\qquad
 \operatorname{Ord}_v(\beta')=\operatorname{Ord}_v(\beta)|_{P_v}.
 \label{eq:canonicalization}
\end{equation}
\end{lemma}
The lemma guarantees that whenever a valid assignment exists, one also
exists that keeps exactly $P_v$ at each node while preserving the relative
order of retained axes.

For example, let $A,B$ have shape $[2]$ and consider the RHS
$\operatorname{sum}(\operatorname{dot}(A,B))$. A valid assignment can keep
the shared axis as batch and sum it later. The canonical assignment contracts
it immediately:
\begin{table}[!htbp]
\caption{Two ways to compute the scalar $A_0B_0+A_1B_1$.}
\label{tab:canonical-example}
\centering
\begin{tabular}{@{}llll@{}}\toprule
Assignment & Dot $(b_\ell,b_r,k_\ell,k_r)$ & Dot shape & Sum axes\\\midrule
$\beta$ & $([0],[0],[],[])$ & $[2]$ & $[0]$\\
$\beta'$ & $([],[],[0],[0])$ & $[]$ & $[]$\\\bottomrule
\end{tabular}
\end{table}
Summation over an empty list leaves its input unchanged.

For TEPL, this means that synthesis can fix the intermediate axis sets
before determining their order and attributes, without losing feasible
rewrites within the defined fragment.

The proof is given in Appendix~\ref{app:canonicalization-proof}.

\subsection{Ordering Feasibility}
\label{sec:ordering-feasibility}
Choosing axis sets is not enough. Each operation also restricts output
order. Let $Req_v$ be an ordered subsequence that the output at $v$ must
contain. Write $A\preceq B$ when $A$ is a subsequence of $B$, and
$A\Vert B$ for concatenation. Propagate requirements from root to inputs
using Table~\ref{tab:order-requirements}.
\begin{table}[!htbp]
\caption{Propagation of required axis orders. A unary node has child $u$;
a contraction has children $\ell,r$.}
\label{tab:order-requirements}
\begin{tabularx}{\linewidth}{@{}llY@{}}\toprule
Location & Required order & Reason\\\midrule
Root & $Req_{\mathrm{root}}=F$ & Match the LHS output order.\\
Permutation & $Req_u=[]$ & The permutation can rearrange axes.\\
Unary summation & $Req_u=Req_v$ & Summation preserves the order of surviving axes.\\
Contraction & \begin{tabular}[t]{@{}l@{}}$Req_\ell=Req_v|_{LF_v}$\\$Req_r=Req_v|_{RF_v}$\end{tabular}
& Free-axis order comes from the corresponding operand.\\\bottomrule
\end{tabularx}
\end{table}
Batch-axis order is chosen at the contraction itself through its attribute
lists.

At each contraction, the required sequence must respect its three output
blocks:
\begin{equation}
 Req_v=(Req_v|_{B_v})\Vert(Req_v|_{LF_v})\Vert(Req_v|_{RF_v}).
 \label{eq:block-order}
\end{equation}
At an input node, its fixed axis order must contain the requirement:
\begin{equation}
 Req_v\preceq I_v.
 \label{eq:input-order}
\end{equation}
Combining these with the root live-set check gives
\begin{equation}
\begin{aligned}
 \operatorname{check}(p)={}&[P_{\mathrm{root}}=\operatorname{set}(F)]\\
 &{}\land\bigwedge_{\text{input }v}[Req_v\preceq I_v]\\
 &{}\land\bigwedge_{\text{contraction }v}
 \left[Req_v=(Req_v|_{B_v})\Vert(Req_v|_{LF_v})\Vert(Req_v|_{RF_v})\right].
\end{aligned}
\label{eq:feasibility}
\end{equation}
These checks depend only on source metadata, which determines the LHS
classes, input orders, and target output order.

In TEPL, the rule's \code{where} condition (Section~\ref{sec:where}) must
imply $\operatorname{check}(p)$ for every legal LHS match within this
fragment. Section~\ref{sec:condition-suggestions} explains how TEPL verifies
this requirement and suggests additional conditions.

\subsection{Constructing Attributes}
\label{sec:formal-construction}
Once $\operatorname{check}(p)$ succeeds, the remaining task is to turn
the selected axes and ordering requirements into RHS attributes.
The $\operatorname{build}$ program processes the RHS from inputs to root,
using each child's output order to determine the required axis positions.

At each node $v$, we choose an output order $\operatorname{Ord}_v$
that contains exactly the axes in $P_v$ and satisfies the required
relative order $Req_v$:
\begin{equation}
 \operatorname{set}(\operatorname{Ord}_v)=P_v,
 \qquad Req_v\preceq\operatorname{Ord}_v.
 \label{eq:order-invariant}
\end{equation}
\paragraph{Input.}
Use the input's existing axis order, $\operatorname{Ord}_v=I_v$.
No attributes are needed.

\paragraph{Permutation.}
Place $Req_v$ first, then append the remaining child axes in their current
order. The permutation attribute lists their positions in the child.

\paragraph{Unary summation.}
Set $\operatorname{Ord}_v=\operatorname{Ord}_u|_{P_v}$. The summation
attribute lists the positions of $P_u\setminus P_v$ in the child.

\paragraph{Contraction.}
Construct the attributes and output order in three steps:
\begin{enumerate}
\item \emph{Choose axis orders.}
Form the batch sequence $b$ by placing $Req_v|_{B_v}$ first, then appending
the remaining batch axes in a fixed order. List the contraction axes
$K_v$ in a fixed order $k$. Use the same sequences for both operands.

\item \emph{Convert axes to operand positions.}
The notation $\operatorname{pos}(S,A)$ denotes the list of zero-based
positions of the axes in $A$ within $S$, following the order in $A$.
For example, $\operatorname{pos}((i,k,j),(j,i))=[2,0]$.
The four attribute lists are
\begin{equation}
\begin{aligned}
 b_\ell&=\operatorname{pos}(\operatorname{Ord}_\ell,b),&
 b_r&=\operatorname{pos}(\operatorname{Ord}_r,b),\\
 k_\ell&=\operatorname{pos}(\operatorname{Ord}_\ell,k),&
 k_r&=\operatorname{pos}(\operatorname{Ord}_r,k).
\end{aligned}
\label{eq:formal-attributes}
\end{equation}
\item \emph{Determine the output order.}
The output lists batch axes first, then left-free axes, then right-free
axes. Free axes keep their order from the corresponding operand:
\begin{equation}
 \operatorname{Ord}_v=b\Vert
 (\operatorname{Ord}_\ell|_{LF_v})\Vert
 (\operatorname{Ord}_r|_{RF_v}).
 \label{eq:constructed-order}
\end{equation}
\end{enumerate}

Together, these attribute assignments form
$\beta=\operatorname{build}(p)$. The rewrite tree is fixed, but the
attribute lists may vary in length, so the same program handles different
input ranks and shapes within the defined fragment. Both $\operatorname{check}$
and $\operatorname{build}$ operate on finite metadata lists.
This is a mathematical program construction; Section~\ref{sec:method}
describes TEPL's symbolic groups and expression generation.
